% Part: first-order-logic
% Chapter: semantics
% Section: intro-semantics

\documentclass[../../../include/open-logic-section]{subfiles}

\begin{document}

\olfileid{fol}{syn}{its}

\olsection{Inleiding}

De semantiek houdt zich bezig met de betekenis van uitdrukkingen. Het
centrale semantische begrip is vervulling in !!a{structure}.
!!^a{structure} legt de betekenis van de bouwstenen van de taal vast:
!!a{domain} is een niet-lege verzameling objecten. De kwantoren worden
zo geïnterpreteerd dat zij over dit domein lopen; aan !!{constant}s
worden elementen van het domein toegekend, aan !!{function}s functies
van het !!{domain} naar zichzelf en aan !!{predicate}s relaties op het
!!{domain}. Samen vormen het !!{domain} en deze toekenningen aan de
basiswoordenschat !!a{structure}. Ook !!^{variable}s kunnen in
!!{formula}s voorkomen. Om daaraan betekenis te geven, kennen we hun
eveneens !!{element}s van het !!{domain} toe---dit heet een
variabelenbedeling. De waarheidsrelatie brengt al deze gegevens ten
slotte samen. !!^a{formula} kan worden vervuld in
!!a{structure}~$\Struct{M}$ ten opzichte van !!a{variable} bedeling~$s$,
genoteerd als $\Sat{M}{!A}[s]$. Ook deze relatie wordt door inductie
naar de structuur van~$!A$ gedefinieerd. Met de waarheidstafels voor
de logische connectieven wordt bijvoorbeeld de vervulling van
$(!A \land !B)$ gedefinieerd aan de hand van het al dan niet vervuld
worden van $!A$ en ~$!B$. Vervolgens blijkt dat de !!{variable}
bedeling niet ter zake doet als de !!{formula}~$!A$ !!a{sentence} is,
dat wil zeggen geen vrije variabelen heeft. We kunnen dan eenvoudig
spreken van !!{sentence}s die al dan niet in !!{structure}s worden
vervuld.

Op grond van de waarheidsrelatie $\Sat{M}{!A}$ voor !!{sentence}s
kunnen we vervolgens de fundamentele semantische begrippen geldigheid,
semantisch gevolg en vervulbaarheid definiëren. !!^a{sentence} is
geldig, $\Entails !A$, als iedere !!{structure} haar vervult. Zij is
een semantisch gevolg van een verzameling !!{sentence}s,
$\Gamma \Entails !A$, als iedere !!{structure} die alle !!{sentence}s
in~$\Gamma$ vervult, ook~$!A$ vervult. Een verzameling !!{sentence}s is
vervulbaar als minstens één !!{structure} al haar !!{sentence}s
tegelijkertijd vervult. Omdat !!{formula}s inductief zijn gedefinieerd
en vervulling op haar beurt door inductie naar de structuur van
!!{formula}s is gedefinieerd, kunnen we met inductie eigenschappen van
onze semantiek bewijzen en verbanden tussen de gedefinieerde
semantische begrippen aantonen.
\end{document}
